% from ega0-22-fr.tex
\oldpage[0\textsubscript{IV-1}]{213}

\begin{remark}[22.7.4]
\label{0.22.7.4-fr}
On a en fait montré (en utilisant (21.9.6)) que la condition b') est encore
équivalente aux autres conditions de (22.7.3) lorsque l'on assujettit $k'$ à
être un des corps d'une famille filtrante décroissante $(k_\alpha)$ de
sous-corps de $K_0$ contenant $K_0^p$, tels que
$[K_0:k_\alpha]<+\infty$ pour tout $\alpha$ et que
$\bigcap_\alpha k_\alpha=K_0^p$.
\end{remark}

\begin{corollary}[22.7.5]
\label{0.22.7.5-fr}
Sous les hypothèses générales de (22.7.3), l'ensemble des idéaux premiers
$\mathfrak n\in\operatorname{Spec}(B)$ tels que $B_{\mathfrak n}$ soit une
$k$-algèbre formellement lisse (pour la topologie $\mathfrak n$-préadique)
est ouvert dans $\operatorname{Spec}(B)$.
\end{corollary}

\begin{proof}
Avec les notations de (22.7.3), il suffit de voir que si $B_{\mathfrak p}$
est formellement lisse sur $k$, il en est de même de $B_{\mathfrak p'}$, pour
tout $\mathfrak p'\in V(\mathfrak q)$ assez voisin de $\mathfrak p$ dans
$\operatorname{Spec}(A)$. Or, si les $k$-dérivations $D_i$ et les éléments
$f_i$ de $\mathfrak q$ vérifient le critère b) de (22.7.3), on a aussi
$f=\det(D_if_j)\notin\mathfrak p'$ pour $\mathfrak p'\in D(f)$~; d'autre
part, il existe un $g\notin\mathfrak p$ tel que les images des $f_i$ dans
$\mathfrak qA_{\mathfrak p'}$ engendrent $\mathfrak qA_{\mathfrak p'}$ pour
$\mathfrak p'\in D(g)$ (Bourbaki, \emph{Alg. comm.}, chap.~II, §~5,
n\textsuperscript{o}~1, prop.~2), ce qui achève de prouver le corollaire.
\end{proof}

\begin{corollary}[22.7.6]
\label{0.22.7.6-fr}
\textnormal{(Nagata).} Soit $B$ un anneau local noethérien complet contenant
un corps. Alors l'ensemble des $\mathfrak n\in\operatorname{Spec}(B)$ tels
que $B_{\mathfrak n}$ soit régulier est ouvert dans $\operatorname{Spec}(B)$.
\end{corollary}

\begin{proof}
En effet, $B$ est une algèbre sur un corps premier $k$, donc (19.8.8, (ii))
de la forme $A/\mathfrak q$, où
$A=K_0[[T_1,\ldots,T_n]]$ est un anneau de séries formelles et $K_0$ une
extension de $k$~; d'autre part, dire que $B_{\mathfrak n}$ est régulier
équivaut à dire que c'est une $k$-algèbre formellement lisse pour la topologie
$\mathfrak n$-préadique (19.6.4). Toutes les conditions d'application de
(22.7.5) sont donc réalisées.
\end{proof}

\begin{remarks}[22.7.7]
\label{0.22.7.7-fr}
\begin{enumerate}
  \item[(i)] Nous verrons plus loin, avec Nagata, en utilisant (22.7.6), que
  l'on peut étendre la conclusion de (22.7.6) au cas d'un anneau local
  noethérien complet quelconque (IV, 6.12.7).

  \item[(ii)] La conclusion du théorème (22.7.3) n'est plus nécessairement
  exacte lorsqu'on ne suppose plus que $[k:k^p]$ soit fini. Prenons par
  exemple $k=\mathbf F_p(X_n)_{n\geqslant0}$,
  $A=k[[T,U]]$, $\mathfrak q$ égal à l'idéal principal $Af$, où
  \[
    f=U^p-\sum_{n=0}^{\infty}X_nT^{np}~;
  \]
  $A$ est un anneau factoriel dans lequel $f$ est un élément extrémal, car il
  est extrémal dans l'anneau de polynômes $k((T))[U]$, donc aussi dans
  l'anneau de séries formelles $k((T))[[U]]$ (Bourbaki,
  \emph{Alg. comm.}, chap.~VII, §~3), et \emph{a fortiori} dans $A$.
  Prenons $\mathfrak p=\mathfrak q$, de sorte que $B_{\mathfrak p}$ est le
  corps désigné par $E$ dans (21.9.7), qui est séparable sur $k$
  (\emph{loc. cit.}), donc une $k$-algèbre formellement lisse. Mais on a ici,
  avec les notations de (22.7.3), $k=K_0=K$, et les $k$-dérivations de $A$
  dans lui-même sont les combinaisons linéaires de $\partial/\partial T$ et
  $\partial/\partial U$~; comme on a
  $\partial f/\partial T=\partial f/\partial U=0$, les critères b) et b') de
  (22.7.3) ne sont pas vérifiés.
\end{enumerate}
\end{remarks}

% (continued in ega0-23-fr.tex, before its first page marker)
\section{Anneaux japonais}
\label{section:0.23-fr}

Les résultats de ce paragraphe seront complétés dans IV, (7.6) et (7.7).

\subsection{Anneaux japonais}
\label{subsection:0.23.1-fr}

\begin{definition}[23.1.1]
\label{0.23.1.1-fr}
On dit qu'un anneau intègre $A$ est un anneau japonais si, pour toute
extension finie $K'$ de son corps des fractions $K$, la fermeture intégrale
$A'$ de $A$ dans $K'$ est

